
\section{On the Free Church’s “Leading Principles”}

Georg Sverdrup

Published in Folketbladet for December 21 and 28, 1898; January 4, 11, 18, and 25; February 8 and 15, 1899.

\subsection{According to God’s Word, the congregation is the proper form of the Kingdom of God on earth}

As I have been urged to write a little by way of elucidation of the Free Church’s principles and rules, I am grateful for the opportunity which Folketbladet has kindly granted me, so that these reflections may reach as many as possible among the people of the Free Church.

And it is also fortunate that we have come so far that, in the midst of struggle and opposition, we can calmly consider our own work without looking to the right or to the left.

A Lutheran Free Church is also something so new that it is not strange if many pause short and think the matter over.

The committee which from the beginning prepared drafts of rules for the Free Church’s work had a very difficult task. It was forbidden to try to form a new faction, at the same time as it was charged with promoting cooperation among free congregations.

But in difficulties of this kind God’s Word shows the way. And for us the only question is to find that “way out” which the New Testament shows, and no other.

For above all other principles in the Free Church stands the decisive authority of God’s Word in the matter of the congregation and in church polity, fully as much as in the question of doctrine.

It is therefore not by cleverness and subtle ingenuity that we try to find the way; it is by childlike fidelity to the Lord’s Word. And therefore there are no living Christians who cannot and ought not to take part in the solution of the congregational question; for the light of the Word shines no less for the simple than for the wise and understanding, and most of all for the former, so that no one is exempt from occupying himself seriously and sincerely with all that belongs to the Kingdom of God. And next to the question of the salvation of the individual soul, there is no more important question than that of the congregation.

“The Kingdom of God is near” is John’s cry, and the preaching of Jesus and His disciples before the outpouring of the Spirit. Since then it may be said that Christian preaching, in this respect, sounds thus: “The Kingdom of God is here.” Yet of course always coming anew and ever reaching farther to more, according to the second petition in the Lord’s Prayer.

The Kingdom of God has its inner essence, which is righteousness, peace, and joy in the Holy Spirit. It also has its outward form, which is the congregation, the Spirit’s people in the world.

It was Jesus’ purpose in His coming into the world to bring the Kingdom of God. And He would not only bring a spiritual benefaction into hearts, but He would also have a new society, to which He Himself has given the name congregation.

What Jesus willed, that He fulfilled by death, resurrection, ascension, and the outpouring of the Spirit. Then the Kingdom of God came in the form of the congregation, with power to live and grow and propagate itself, so that “the new creation” might become manifold and fill the earth as the old did; but by new means, God’s Word and Spirit, this enlargement and extension was to take place.

For this form of the Kingdom of God, which is the congregation, all the apostles and all the first Christians thereafter labored, because the Spirit drove them to it. And within the revelation of the New Testament there is no mention of any other form for the Kingdom of God.

We hold that in the New Testament there is no mention of any bishopric over, or within, more than one congregation, nor of any papacy, nor of any church department, nor of any church council or council or synod. There is a congregation in every place where there are Christians, and this congregation has its elders or bishops; but there is no “church government” of any kind.

It is a sign of decline and decay that, not long after the days of the apostles, men began to get “church governments” of one kind or another over greater or smaller parts of the church, and thereby thought to obtain a better form of the Kingdom of God. In the actual history of the church there has never at any time been any common church government over the whole Christian church. Nor is there any such thing now in our own day.

It therefore seems to us an unreasonable and unjust demand when it is now often presented in such a way that all congregations are obliged to belong to this or that outward association, so that they may come under the proper church government. Those are called anarchists who do not regard a common church government as necessary and as a compulsory matter. But what was not found in the apostolic age, what has never since existed in the church, and what does not exist in the present time, can scarcely be necessary to Christianity. For if this were necessary, then there has hitherto never been the right kind of Christianity.

For the Kingdom of God it is, according to the New Testament, necessary to have congregation, but we cannot see that any other outward organization over the congregation is a necessary part of Christianity.

It is therefore a principle in the Free Church, which bows only to God’s Word and to nothing else, that the congregation is a sufficient form of the Kingdom of God, and that no other form is required from the outpouring of the Spirit until Christ’s return. But according to God’s will, revealed in His Word, the congregation is in return absolutely necessary where it is possible. Where a single Christian lives alone in prison or under some other constraint, the lack of a congregation will not become for him unto perdition; but where Christians can be a congregation and will not, there is surely great danger for their souls.

As to whether councils, synods, popes, and church departments, which are not commanded in God’s Word, have done more good or harm, we cannot judge; God knows.

Among us here in America these authorities over the congregations have in any case contributed much to factionalism and other sin.

Nothing, therefore, can be more timely and more required among us than a work that aims at the restoration of the congregation in biblical, and especially in New Testament, form. And there is nothing that works more discouragingly upon a Christian mind than to see that so many Lutheran pastors oppose this cause with great zeal.

If it is good for Christian people and for the Kingdom of God on earth that congregation be gathered, and if it is God’s Spirit that gathers congregation, why does one not grant the people this good?

In many places the only question is whether the people want a pastor. When it became a matter of serious inquiry whether they also wanted congregation, then the opposition also began in earnest from the side of many pastors. But if the congregation is a great and glorious divine blessing, for the sake of which Christ has come and the Spirit has been poured out, why does one not want human beings to become partakers in it?



% Page 96


\subsection{The congregation consists of believing human beings, who by the use of the means of grace and the gifts of grace according to the direction of God’s Word seek their own salvation and eternal blessedness, and that of all their fellow human beings.}

The congregation is the fulfillment of Moses’ sigh and prayer in Num. 11:29: “Would that all the Lord’s people were prophets, would that the Lord would put his Spirit upon them!”

The congregation is the fulfillment of David’s song in Psalm 110:3: “Your people come willingly in the day of your power; in holy splendor your youth come to you like dew from the womb of the dawn.”

The congregation is the fulfillment of Jeremiah’s prophecy in Jer. 31:33: “But this is the covenant that I will make with the house of Israel after those days, says the Lord: I will put my Law within them and write it in their hearts, and I will be their God, and they shall be my people; and they shall no more teach each his neighbor and each his brother, saying, Know the Lord; for they shall all know me, both small and great, says the Lord; for I will forgive their iniquity and remember their sin no more.”

The congregation is the fulfillment of Joel’s prophecy in Joel 3:1: “And it shall come to pass afterward that I will pour out my Spirit upon all flesh, and your sons and your daughters shall prophesy; your old shall dream dreams, your young men shall see visions; and also upon the menservants and upon the maidservants in those days will I pour out my Spirit.”

And a multitude of other testimonies from the Old Testament might be brought forward, which bear witness that the Lord in the fullness of time would set a new people of God upon the earth, who should not suffer under the deficiencies under which the old Israel suffered.

What the prophets always complained of in the old Israel was this, that it answered so poorly to its calling and its election, because there were so few who were truly filled with the Lord’s Spirit and had living faith. Most had sunk into spiritual lifelessness and formalism, not to mention that many even lived outright in sheer worldliness and gross idolatry.

But the Lord promises that the sighing of his children shall be heard and their complaint and lack remedied, and the day shall dawn, a “morning without clouds,” when a new people of God shall be born, like the dew from the womb of the dawn, in which all shall have the Spirit, and all shall have forgiveness of sins, and all shall have the Law written in their hearts, and all shall know the Lord, and all shall prophesy.

These are indeed the promises. It is against these that one sins altogether, when one teaches that two or three Christians in each assembly are enough to make a congregation out of it.

Jesus Christ himself says: “You are Peter, and upon this rock I will build my congregation.” And already from that one ought to understand that congregation is not merely a place where God’s Word is preached and the Sacraments administered, or a district with geographical boundaries, or a gathering of hearers of God’s Word. No, Christ’s congregation consists precisely of the same kind of Spirit-filled persons as Simon son of Jonah, whose knowledge and confession were not of flesh and blood, but from the Father in heaven.

And what the prophets had foretold, and Jesus promised, that God did when he sent the Spirit at Pentecost after Jesus’ resurrection and ascension. Then the congregation was born; and the signs that accompanied the birth are precisely the signs of the Spirit and of life, in order that it may be established for all times that where there is not the mighty power of the Spirit, and where there is not the fire of love, and where there is not united witness to the Lord’s great work, there is no congregation.

This newborn congregation was a visible company of human beings. Not only the Word and the Sacraments were visible; the persons who constituted the congregation were visible. The sound of the mighty storm-wind was heard out through the city, so that the multitude came together; the tongues of fire which settled upon each of them showed who it was the Lord had filled with the Spirit; the witness to the Lord’s great works was heard, and those speaking were seen by all, both believers and unbelievers.

So far the congregation is visible; and when 3,000 were added to the congregation by baptism, then the congregation was still visible. And when all those who believed were steadfast in the apostles’ teaching and the fellowship and the breaking of bread and the prayers, and when they were of one accord daily in the temple, and when they were together and had all things in common, then it was a visible congregation of believing human beings.

Was all this then a mistake, a fanaticism without meaning and without significance for the future? Was this a misunderstanding which history was to correct? Or was it the light, the light of the Spirit and of truth, which was to shine for souls and for congregations through all times until Christ’s return?

We hold that this was the truth and is the truth and remains the truth; and that what is called congregation is bound to examine itself by this light.

One says that it is meaningless, for all know that we cannot rise to this pattern. But is not Jesus Christ a pattern for us, and are we not bound to walk as he walked? We know indeed that we cannot attain it; but does it not become us then to live in acknowledgment of sin and humility? And for the congregation even this would already be an advance, if it would begin to live in repentance instead of in carnal security and self-assurance. Hold up the light, and be assured that it never becomes too bright in the Church.

But is it impossible? Is it impossible for us to repent and believe in Jesus Christ? Let us not give up hope for the conversion of any human being or any member of the congregation. So long as there is life, there is hope. The Lord will pour out his Spirit upon all flesh; the question is whether all flesh is ready to receive the Spirit; for Spirit cannot be forced upon anyone.

Thus according to the New Testament there can be no doubt about the essence of the congregation. It is the people of the Spirit, so that it extends exactly as far as the working of God’s Spirit. It is Christ’s body, it is Jesus’ bride, it is God’s living temple of living stones, it is God’s household and his dwelling in the Spirit.

One will say, yes, that is the whole Church; it is so good, but each little congregation is so poor.

Not at all. That first congregation in Jerusalem was indeed both the whole Christian Church and the single local congregation. On that day the one was no greater or wider in extent than the other. And nothing has ever become good by becoming large. The congregation is the “little Church,” as the learned are wont to say. And by that is meant that all that the Church has, the congregation also has; and all that the Church is, the congregation also is. And besides, it is surely clear to all thinking people that each single little part of a building can more easily be good and strong and complete than the whole building; for among 100 stones a bad stone may easily be found, while perhaps the 99 may be good stones; and while the one bad stone does not alter the 99 good stones each in itself, it can nevertheless do much harm to the whole building. In any case it is clear that the responsibility falls upon each single congregation; it is that which first and foremost is to be God’s house and the body of Jesus Christ.

But if the congregation consists of believing human beings, is it thereby meant that it is perfect, that is, without sin?

No, by no means. The Christian persons are not perfect, but shall become so. Or was it not said to Peter—perhaps the very same day he was called the rock—Get behind me, Satan! There is no perfect congregation before the resurrection of the dead.

Let us await him who is to come again to fetch the bride to the wedding of the Lamb, and let not the bride fail to prepare herself. And while we wait, let this be our daily prayer: Come, Lord Jesus, come soon! For where this expectation stiffens, whether because one has sunk back into the world, or because one thinks one has already come to the goal, there the congregation is no longer living.

The congregation has means of grace and gifts of grace, with and by which it works. The common explanation, that the congregation has means of grace, and that it has the power and authority to choose whom it will, believer or unbeliever, to administer them, only provided he “has passed his examination,” is not Christian, even if it suits well our state-church institutions and usages. The congregation also has gifts of grace, if it is faithful in its calling. Indeed every believer has his gift of grace and is bound to use it for the edification of Christ’s body.

For all means and all gifts have this one goal, which is the same as the salvation of souls, or the spread and growth of the Kingdom at home and abroad. We are all called to work in the vineyard, each according to his ability and gift, each with his pound. And whatever our station in life may be, if we are in the congregation, then this is our calling: to live for God and for his Kingdom, to labor for it with all the zeal that God’s Spirit works and Christ’s love creates.

For if the congregation is the people of God’s Spirit, then it is the instrument of God’s love upon earth. But God’s love is in a wholly particular sense that which seeks after the lost in order to save it. For many reasons it is necessary to lay great weight upon this last point just now. For the strayed and scattered sheep are surely not far from any one of us. And this is the most important of all signs of life in the congregation, that it wills to have all the lost saved. All those here at home, all those out there, all in the Christian lands, all heathen who do not know God. World-embracing is the congregation’s calling, and yet the work for its fulfillment is so near at hand that it sounds to every soul: Go today into my vineyard!

If the believers in each outward congregation will begin this work, then there will soon be change among us, and in the midst of our spiritless age it will become manifest that God’s Spirit is still mighty to bring forth congregation.

% Page 102



\subsection{According to the New Testament, the congregation needs an outward organization, with a register of its members, election of officers, fixed times and places for its assemblies, etc.}

The congregation, which from the beginning is born of the Spirit, can live only by the Spirit. Therefore the Spirit’s means, the Word and the Sacraments, are the means by which the life of the congregation is propagated and preserved. It is only Baptism, the Word, and the Lord’s Supper that can bring forth a living congregation and sustain it, when these means are rightly used.

It appears that some have come upon the thought that it is the organization and the constitution that decide the matter. And perhaps it is not so few who have this thought, though in somewhat different fashion. Some think that if a person only enters the orthodox congregation, then he is a Christian; for the congregation, after all, consists of Christians. Others think that if a person enters our congregation and is acknowledged by us as a Christian, then it is—if not certain—yet at least probable, that he is a Christian. Both sides—as opposite as they are to one another—perhaps lay greater weight upon the organization and the constitution than is right. For neither can a constitution or organization make alive, nor can it sustain life.

Nevertheless, we maintain with full conviction that the congregation needs an outward organization, and that this need is already known and met in the New Testament: the Lord did not conclude the revelation of and concerning the congregation before the need had shown itself and had been remedied.

We lay no little weight upon this matter, since there are not so few voices to be heard which reject and almost despise every outward organization of the congregation.

The necessity of this outward organization is not the same kind of necessity as that of grace and of the Word and of the Sacraments. That is to say, it does not serve directly for the creation and sustaining of life. But surely we can say that it serves directly for the preservation and protection of life.

In other words, organization does nearly the same service for the congregation that good clothing and a good house do for a man or a family. And this service is by no means so very small under most circumstances.

As the body exists before the clothes, so also the congregation exists in the New Testament before the outward organization. It is the congregation that itself organizes itself by counting and recording its members, by choosing its officers, by determining when and where one shall assemble, and how things shall proceed in the assemblies. These things, and perhaps others that someone might mention, are clearly and plainly set forth in the New Testament; and without doubt Christian people do well to heed them, for he who will not heed the will of God also in such matters cannot expect divine blessing in full measure.

That the congregation may count its members and keep them in a register follows, according to the testimony of the New Testament, quite straightforwardly from the fact that it receives its members by Baptism. To baptize is in its form such an outward act that it is very simple to count how many are baptized. That none was compelled by civil law to be baptized, and that none was baptized who did not confess his sins and his faith in the risen Jesus, is surely self-evident to every attentive reader of the Bible. But, as said, this outward act made the persons easy to know and to count.

And that the congregation had to choose officers and make other determinations—which arose rather from the congregation’s imperfection than from its perfection—involved the use of the vote, which according to Acts 6:2 and other places belonged to “the whole multitude of the disciples,” so that there was no difference of rank between governors and governed; they were all sharers in the government, if they were members of the congregation.

There is also no doubt that from the beginning the congregation had its regular assemblies, and that these assemblies, in all their unique freedom, were nevertheless beautifully and decently ordered. Read Paul’s movingly simple description of the assembly in 1 Corinthians 14:26: “What then, brothers, is to be done? When you come together, each one has a psalm, he has a teaching, he has a tongue, he has a revelation, he has an interpretation; let all be done for edification.” And the following verses in the same chapter give still more exact insight into the life and the order that were found in the assemblies. There is no stiffness and no excited vehemence; there is also no disorder, and no arbitrary self-exaltation of one at another’s expense.

Surely, the Lord has not left His congregation to human cleverness and subtlety, as many think, so that we have no instruction concerning “church polity” in the Bible, but must speculate our way to something “suited to the times”; there is light enough for the congregation in the New Testament, and it is not worth while to let anyone persuade himself that it does not fit “in our time.” It did not suit worldly people any better then than it does now. It always has been and will continue to be foolishness to the Greeks and an offense to the Jews, together with the crucified Christ and His congregation.

And here there is only one thing that we will point to in conclusion, and which perhaps some will think a little about. Every Christian congregation ought to be fully self-sustaining in the question of its own edification. For three or four congregations to have one pastor in common is a rather imperfect arrangement; and it becomes so defective that it is almost intolerable when there can be no service because the pastor does not come. When the congregation does not even assemble every Sunday, how then can it be expected that life and fellowship should thrive? If the congregation cannot assemble for edification when the pastor does not come, then every congregation ought to have its own pastor, and it should surely prove that this was not too much. And if a congregation has become so far Christian that it can build itself up with God’s Word even when the pastor is ill or in the neighboring congregation, then with a little experience it will surely discover that the more living the congregation is, the more it will need the pastor’s labor.

I know that some will answer that it costs too much, and to that I will say nothing more than this, that the first congregations were scarcely so rich as is commonly the case among us; yet they met their expenses chiefly on account of their strong organization. At this point too, some of the clever heads of our time could learn much by studying the New Testament and its instruction for ordering the financial affairs of the congregation. And there is no doubt that in this matter also better results would be attained by strict accordance with Scripture than by all sorts of worldly contrivances and arrangements.

If our people would bestow as much diligence upon the congregation’s organization as upon the so-called societies, then the whole aspect of churchly work would be quite different from what it now is.

Whether one dissolves the congregation by laying all weight upon the church association, or one dissolves the congregation by laying all weight upon the individual person, one does equal harm and acts equally wrongly. In the final result, for the dissolution of the congregation, there is no difference between the stiffest high-churchmen and the loosest “free-free.” Both parties, in truth, abolish the congregation.

Hold fast to the congregation, and take pains with its exact and careful organization! Only a well-organized and united congregation is fully capable of labor. For organization is in truth nothing other than every man and every woman at his or her post and in his or her work.




% Page 107


\subsection{Since not all members of the congregation’s outward organization are always believing people, and since such hypocrites often take false comfort from their outward connection with the congregation, it is the congregation’s holy calling to cleanse itself both by awakening proclamation of the Word of God, by earnest admonition and reminder, and by exclusion of open and obstinate sinners.}

The congregation consists in reality of believing people. And just as every human being is instructed by the Word of God that he cannot be saved and blessed except through faith in Jesus, so all are likewise instructed by the Word of God that only when they are true believers are they also actual members of the congregation.

When, however, there is question of the congregation’s outward organization, then the Word of God and Christian experience down through church history bear witness that it both can and will happen that such persons are included in the organization as do not in spirit and truth belong to the congregation. Likewise, that it can and will happen that some, who began in the Spirit, end in the flesh, or fall away from the faith in their hearts, even though they retain the outward connection with the congregation’s organization. For these things take place in the heart’s secret hiddenness, and there are no definite outward signs of them.

This fact has its simple explanation in this, that the congregation’s outward organization must take account of the outward signs and marks of the congregation or of Christian life, and cannot judge the hidden state of the heart.

We say that the Word of God bears witness to this fact. And then we think not so much of the seven parables mentioned, as of the Acts of the Apostles, the Epistles, and the Revelation of John.

For of Jesus’ parables one may indeed rightly say that they speak rather of the effect of the Word in general than specifically of the congregation’s organization. For example, when the kingdom is compared with a dragnet that catches good fish and rotten fish, then the net is rather the Gospel in its work upon human souls than it is the congregation’s outward organization. Yet it is also right to say that at least several of the parables can be applied to the congregation’s organization, as for example the workers in the vineyard and the ten virgins.

But when we read the Acts of the Apostles and the Epistles and the Revelation, then we receive a full and clear light upon this matter, that the organized congregations, according to their outward boundaries, comprise not only vessels of gold and silver, but also of wood and clay, and some for honor, others for dishonor. It is also clearly shown that the New Testament contains the truth as it is applied to the most various spiritual conditions, and it is presupposed that these are found within the organized congregations.

In other words, just as each individual Christian is never all that he ought to be, so it is also with the congregation, and admonition, warning, encouragement, comfort and consolation, instruction and teaching must be used within the congregation’s outward organization, in order that the sleeping may awake, the wavering be kept from falling, the doubting be established in the truth, and so forth. As a rule, the whole Word of God will find its application within the congregation; at any rate the apostolic Epistles, with their varying content and their many exhortations, were written to Christian congregations.

It almost awakens astonishment, and it ought to awaken horror, that it was so already in the first congregations. Yet God had in a quite particular way governed all things so that the congregation’s outward organization from the beginning should comprise believing people and no others.

What was John the Baptist’s work other than to prepare for the Lord a people made ready, that is, to prepare the congregation by preaching repentance and baptism? Within the people of the old covenant there was by John’s ministry to take place a separation of those who would truly flee from the judgment and receive salvation. And John set as the separating mark the baptism of repentance, with which he baptized those who confessed their sins. And after John had thus worked in Israel in a separating and setting-apart manner, Jesus came, and He did not establish the congregation while the people by the thousands thronged around Him and wanted to make Him king; He waited until persecution and crucifixion and death had frightened away all who had carnal and worldly notions concerning Him and the kingdom. And only when the flock had become so little that it seemed wholly useless as a beginning to the kingdom, only then did Jesus make this little flock into a congregation by sending it His Spirit.

Truly, there is something solemn in this that God has done in order to prepare and bring forth the congregation. And those are not right who teach that this was mere exaggeration and fanaticism, and who think that they know much better what the congregation is and ought to be than John and Jesus and the Spirit of truth Himself.

And yet unworthy, hypocritical people pressed their way into the congregation’s outward organization. And while this fact at the same time speaks loudly of the power of divine truth, in that it shows that man gladly wants to appear to be on the right way, even when he is not, this fact must nevertheless strike us with horror at human falsehood and wickedness, that they will sneak in by dishonest means there where only the highest degree of sincerity and honest acknowledgment of sin gives rightful access.

It is therefore not to be wondered at if in the so-called congregations of our own time also we find many false members. But it is in truth both to be wondered at and to be horrified over, that there are ministers who teach that this is right and proper.

When for centuries one has by the power of the Law, the ban, and the national church forced people into the church (we will not even call such a thing a congregation), then of course it is self-evident that matters must be in a bad state. And matters are in a bad state. Not merely because gross vices are found in the midst of the great majority of congregations, but most of all because one makes no attempt whatever to return to the congregation’s concept and essence.

What then is to be done about it? Fold the hands in the lap and comfort ourselves that we are, after all, not Sodom and Gomorrah?

Is that any comfort? Or is it not written: It shall be more tolerable for Sodom and Gomorrah on the Day of Judgment than for you? Where then is the comfort?

No, what is to be done? First and foremost the same as John the Baptist and Jesus and the apostles did. That is the main thing, and all else is secondary, though it belongs to the work as well.

And we all know what these did; but there is so frightfully little done in following after them.

They preached the Word without favor and without fear; wholly and truly, personally and directly, without respect of persons.

And this is and remains the main thing: light, light, and yet more light in the church. It never becomes too bright, nor bright enough. And upon this the congregation’s true thriving depends. In the light from God, the good thrives, just as surely as the evil fails to thrive. In God’s truth there is nourishment for the souls that seek the truth; in God’s truth there is also cutting judgment over all dishonest and hypocritical hearts. Let the truth sound forth in full, and there shall come both brokenness and resistance, both faith and hardening, both love and hatred.

But the truth is Jesus Christ. Only by a real preaching of Christ does a congregation come into being. Let that sound forth, and the congregation is built and grows and is strengthened; put something else in its place, however popular it may be, and the congregation sinks into a swamp and is corrupted and loses more and more of its true essence and its true strength.

This is the struggle for life in the congregation. It must be kept going all the time, as long as the congregation is here below. It never stops, except where spiritual death takes complete control.

And God, who is rich in mercy, has graciously led and ordered it so that even where everything seems to be death and sleep, He has raised up or sent instruments that have kindled the fire again where it had been extinguished.

But where the proclamation is as it should be, there comes another means in addition, which we also will mention: work.

The congregation lives and exists in order to work. It is a gathering of workers. The congregation is not organized unless there is work for all. It is not a congregation where a gathering of people “hire” a minister to work for them, so that they may be spared it themselves. No, the congregation may in that respect be likened to a choir, which must have a choirmaster to lead and instruct it. Shall the instructor sing alone? Shall not all take part? So it is in the congregation; it is a gathering of God’s servants, men and women, who labor for God’s kingdom, for Christ’s cause. The leader does not become superfluous, but all the more necessary in such a company.

But this work will also help to raise up the congregation and keep it free from worldliness. If all have enough to do with spiritual work and spiritual interests, then worldly elements will surely find the congregation quite a narrow place.

Let the work become so heavy to flesh and blood that it becomes a real cross to bear in Jesus’ footsteps, and it shall help the congregation exceedingly much to be truly a congregation. Then the point will not be to become many in order to make it easy, but the point will be to become living strong, in order to do the Lord’s work with power.

The congregation which will not by these means resist its own becoming worldly will soon lie under, and it will become a disgrace to God’s name instead of unto honor for the Lord and unto light for the world.

When therefore we acknowledge that the congregation must have an outward organization, and at the same time acknowledge that there is occasion for hypocrites to sneak into the organization, then this is the strongest exhortation to be on guard and to use the means the Lord has given us, so that the congregation’s enemies may not become its masters through the congregation’s sloth and unwillingness to take up its cross and bear its burdens. Only the Gospel of the cross and the cross of the Gospel can preserve the congregation, so that it does not become a shameful distortion of Christ’s body and Jesus’ bride.


% Page 113


\subsection{The congregation governs its own affairs under the authority of God’s Word and Spirit, and acknowledges no other ecclesiastical authority or tutelage over itself.}

\subsection{The free congregation esteems and loves all the spiritual gifts which the Lord bestows for its edification, and seeks to quicken the gifts and promote their use.}

If it is true, as God’s Word bears witness, that the congregation consists of the believers in each place who there on that place have formed an organization in order to labor for God’s kingdom, then it follows of necessity therefrom that the congregation is self-governing and independent of other men in this its work.

For Christianity is in its innermost essence a relation to God, which brings with it a complete dependence on Him and yet a complete freedom in Him.

If therefore other men have no authority over the congregation in its congregational work, then in the fullest possible measure the congregation does have an authority over it, namely the Word and the Spirit. And so absolute is this authority, that the congregation is no longer a congregation when it tears itself loose from this; and only so far is the congregation a congregation as it follows this authority.

It is this which gives the congregation its self-dependence and its responsibility. It acts as it acts because it is fully persuaded in its own mind that thus it is right and proper according to God’s Word, and thus God’s Spirit drives it through the Word. Therefore it knows also that God will call it to account for every single one of its actions.

For God has not left His believing congregation without light and without strength.

Over the way and work of God’s congregation the light of the Word shines clearly and mightily; and God’s Spirit gives ability and strength to walk the way and to carry out the work.

It is not seldom heard in this connection that the gifts of grace are so scanty nowadays that one properly speaking finds them only among the clergy. And among the clergy they are found, not because God’s Spirit has given them gifts of grace, but because the congregation has chosen and appointed them to the “office.”

This false doctrine is widely spread also among us and is frequently used in defense of the ungodly opinion that an unconverted pastor is a true pastor, provided only that a congregation has called him in an outwardly lawful manner.

A more forcible distortion of God’s truth is scarcely to be found. And yet it is not rare to hear such opinions advanced and defended within the Church.

It is true that no one can be a true pastor, such as a pastor ought to be, unless he has both true faith and the Spirit’s gift of grace for the proclamation of God’s Word; but it is also true that he does not receive the gift of grace, nor faith either, by receiving through the congregation’s election the outward position as preacher or pastor. It is also true that he is not the only one who has received a gift of grace. To all His believers the Lord has given gift and ability to labor for the cause of God’s kingdom.

It is likewise not true that all pastors have gift and power to preach and govern, while all laypeople have only gift and ability to support the office by paying.

The order according to God’s Word is altogether different. According to Scripture, the matter stands thus: the Lord, who ascended on high after the humiliation, crowned His victory by distributing gifts among men. And these gifts made some fit for one kind of work in the congregation, others for another. If then the congregation is to choose its pastors and officeholders in a right and scriptural manner, it must try to choose those whom God has made fit, and not others. Therefore it is written concerning one of the first elections—perhaps the first actual congregational election: “Therefore, brothers, select from among you seven men of good report, full of the Holy Spirit and wisdom.” There a rule is laid down for a right election, and only when the congregation understands and acknowledges this rule and acts according to it will its election be pleasing to the Lord and profitable to the congregation.

It is a great demand upon the congregation’s spirit and insight that it should know those who not only have good report, but also are “full of the Holy Spirit and wisdom”; and this agrees poorly with all the theologians’ frivolous, superficial, and thoughtless use of the saying that “we are not knowers of hearts.” According to Scripture, the congregation is bound to know who are “full of the Holy Spirit and wisdom,” if it is truly to be a free and self-governing congregation.

Thus then the congregation receives its true officeholders, by the Lord’s gifts of grace and by the congregation’s upright election. And it may surely be said, without too great an error, that the congregation’s officeholders are in a special sense the congregation’s workers, cut out for that by God and by the congregation. But by this it is not said that only these shall labor and no one else. All believers, great and small, young and old, men and women, shall be workers in the vineyard. There is no sense at all in understanding the matter so that the members of the congregation lay their work down upon the officeholders. These receive their separate charge, but there always remains more than enough over for all the others.

Take, for example, the Christian upbringing and instruction of children. Can the school of religion abolish parents’ duty and obligation to bring up their children in the discipline and admonition of the Lord? And so in all the spiritual branches of congregational work. The more spiritual a matter is, the more are all believers bound to take part in it with their gifts and abilities.

Naturally, we understand well the significance of outward order, so that when the congregation elects a treasurer, all the others cannot be treasurers at the same time. But they all can and shall help the treasurer to keep the books in good order. Yet just as idiotic as it would be to suppose that the treasurer alone is to procure the money and manage the money, and no one else is to have the very least concern with that matter, quite as idiotic and a little worse besides is the common conception that all edifying work in the congregation is a matter entrusted entirely to the pastor, and that no other Christian either can or shall contribute anything to that matter with his spiritual gift or ability.

But the congregation must not only use the gifts it receives; it also lies upon it to quicken and cultivate, to draw forth and call forth the gifts.

Here first and foremost the public proclamation of the Word must be called to account. It must not only be said to a man and to a member of the congregation’s outward organization that you must give the Lord your heart in a true and living faith; but also that you must present yourself as a living, holy, and God-pleasing sacrifice through a true and inward and self-sacrificing love to God and to your neighbor. The proclamation must become congregational, so that light falls upon the Christians’ calling and task as members of the congregation, upon their entrusted pounds and talents.

And not that alone. Brotherly spirit and mutual confidence must reign within the congregation, so that one converses and deliberates with one another precisely about these questions: What shall we do, and how shall we do it? Which of us is fit, and what is the duty of the one who is fit?

If one never seeks after the gifts of grace, one surely does not find them; then it sometimes comes to pass that they almost must thrust themselves forward; and that is one of the reasons why many gifts are among us partly wasted and partly warped.

We cannot afford this. The congregation ought to search after the gifts as after noble and precious gold, and then “test everything and hold fast to what is good” in this respect also.

This then is the congregation’s true form in the world: a people of believing, praying, and working men and women, who receive from the Lord all that they need in their work, both the faith and the life and the gifts, and who therefore acknowledge this wholly and fully, that they belong to the Lord with all that they are and all that they have, so that they all are the Lord’s servants, male and female, created by Him for every good work.

Are we such congregations? If not, why not?

And if anyone says that such a thing is fanaticism, that it is altogether impossible, then we answer: Have faith in God; what is impossible for men, that is possible for God!

VI.

\subsection{The free congregation receives with joy the mutual help which congregations can render one another in the work of advancing God’s kingdom.}

\subsection{Such mutual help consists partly in this, that the spiritual forces of the one congregation come to the benefit of the other by free meetings, mutual visits, lay activity, etc.; partly in this, that the congregations voluntarily and at the prompting of God’s Spirit cooperate for the solving of such tasks as may surpass the strength of a single congregation.}

\subsection{Among such purposes are named in particular a pastors’ school, the spreading of Bibles and other writings and periodicals, home mission, heathen mission, deaconess homes, children’s homes, and other works of mercy.}

If the Free Church in its leading principles lays the chief weight upon the congregation and its spirit, life, and freedom, it is very far from being opposed to cooperation among congregations.

We do not deny that the old organizations of synods and councils have taken a fair name when in Norwegian they often call themselves fellowships. We have only a painful experience that at least a part of them do not answer to the name.

Instead of “fellowship” being the free cooperation and the mutual influence of spiritual forces, we have among us party organizations, where that which does not go in the leading-strings of the church-political “machine” is shut out and cut off, and one is as though afraid of a breath of freedom, so that one hates everything that is free, even to the very name of it.

That large parts of our people by means of these organizations have been led into the fanaticism of party-spirit, and into the bitter, spiritless strife and contention that follow from it, is something which all nobler spirits among us agree in lamenting. But many of them are still of the opinion that this frightful form which the “Church” presents to the people can be remedied by “new fellowships” and constantly “new fellowships,” which in reality are little more than new parties with new slogans and new fanaticism and new bitterness.

The Free Church’s guiding principles seek to show a new way by the abolition of parties.

It is not a question of abolishing and doing away with fellowships, but of abolishing party-spirit in the fellowships.

Although therefore all church politicians know well the rapidly kindling fire of party-spirit and the immediate power that follows from it, the Free Church has nevertheless in its principles declared that it renounces this false power, so that freedom may have a little room among us. Perhaps the day may dawn when there can be a cooperation of Christians and of Christian congregations which is led by the Spirit, not by church politics.

We see well enough that the goal is high and distant; but if only the congregation first becomes congregation, then the union of congregations too will surely become a real fellowship.

Meanwhile we must try to practice whatever little we are able of free and spiritual cooperation and of mutual spiritual influence.

Therefore we believe in free meetings, where men and women from different congregations come together to exchange thoughts concerning what is serviceable for the advancement of God’s kingdom. And the freer these meetings are, and the more they break through and strike down the thorn-hedges which party-spirit has set between our people, so much the better.

Let them come together! Those from the one and those from the other “fellowship”! Those who by “fellowship” have lost the real fellowship, and those who without “fellowship” have found the unity of the Spirit in the bond of peace! Provided only that we stand on New Testament and Lutheran ground, we may well meet together and speak with one another concerning the awakening and preservation of spiritual life, concerning the enlivening and liberation of our Church. This is a consideration that is equally useful for the one party as for the other, and these are matters that can be practiced without regard to “fellowship.”

Our Norwegian Lutheran Church with its stiff, spiked, and jagged fences, with churches that pastor shuts against pastor and congregation against congregation, although they have the same doctrine and confession, presents a pitiful and lamentable sight.

Something must be done in order to come out of this shame and wretchedness. While the parties labor to preserve the evil, the Free Church must try by sober, brotherly, and persevering labor to remedy and overcome it.

To this end mutual visits, especially by laymen, will serve greatly, if they are carried out in the right spirit. The fellowship agents cannot do the work that is needed for this aim. The lay preachers ought to understand that it is not for the party they are to labor; it is for the Spirit and the life. We may perhaps come so far one day in the Lutheran Church that even a pastor may preach in another pastor’s church unhindered by the fences; but for the present the lay preachers must make serious business of an impartial awakening work, and the conviction must press through more and more that it is sin to use the gifts of grace in the service of party-spirit, and that it is sin to shut the churches against the proclamation of God’s Word.

Let the very greatest leaders of the fellowships come forward among us and say whether even they must not admit that such a program is justified and such a work required.

But, one says, what will become of such common works as the congregations must do? What will become of mission and school and children’s homes?

The question in truth answers itself. It will be no otherwise than it already is for many such works. Those who were most interested in the mission to the Jews began to carry it on and to work for it. They received gifts and contributions from whoever was willing, and there was no compulsion. Has it not gone well enough? We think, has it not gone better than if it had been a single party that had managed it?

That those most interested should be allowed to tend a cause is freedom’s demand and the gain and advantage of the cause itself. Why should it not be able to go?

Those who have church-political tendencies and would exalt themselves by help of the great force of ecclesiastical interests do not like this. But in return the interests thrive best when not too many politicians are sitting on the load.

This seems to us to be the advantage of the free work. It gathers the interested, and that is fair enough. Those who are not interested ought not either by the fellowship-machine to be compelled and driven to take part in that which lies far from them.

Yes, one says, but many of these interested persons then come, as it were, to appoint themselves to the work; they ought to wait until someone elected them. But suppose now that God put into their mind what they should do, and they followed His call to heavy and burdensome labor?

Or is [1–3 words unclear] less a missionary because he went himself to the Santals and no one sent him? Must God also wait for fellowship resolutions?

Where the Lord’s Spirit is, there is freedom, freedom to labor, freedom to suffer, freedom to live for a cause, freedom even to die for the cause. Should it not be possible to believe also in the Holy Spirit?



% Page 122



\subsection{Free Congregations have no Right to demand that other Congregations submit to their Opinion, Will, Judgment, or Determination; therefore every dominion of a majority of Congregations over a minority is to be rejected.}

\subsection{The common organs that may be found desirable for the Congregations’ cooperation, such as greater or smaller meetings, committees, officers, etc., cannot in a Lutheran Free Church impose upon the single Congregation any duty, bond, compulsion, or burden, but are only entitled to come forward with proposals and exhortations to Congregations and persons.}

In civil society it is long since recognized that pure and absolute majority rule is of dubious worth. Modern society ought properly to consist in the employment of all abilities and powers for the common good; but then it is plainly the less fitting that one should, for example, in a society of 100 exclude 49 from all influence because they are not so many as 51. For experience shows that the 49 may be just as capable and as well-intentioned as the 51. People are not measured by the ell-measure, neither can one in all cases attain the best results merely by counting. It is enough that, in most civil societies—and thus also among us—provision has in one way or another been made for a counterweight against simple and absolute majority rule.

In the Church this question is infinitely more serious, more difficult, and more significant than in the State. And for many different reasons it is an altogether different thing in the Church than in the State to apply a wholly unrestricted majority rule.

In the United Church they have tried to carry it to an extreme, which I believe was done not from principle, but from considerations of expediency. It was a convenient means of displacing Augsburgers, who were in the minority. And because it seemed convenient, they therefore went to such an extreme as is impossible in civil society; they not only denied the minority all influence, even the most rightful, but simply drove it out of the body and then called this expulsion—church discipline. That, thanks be to God, one cannot do in the United States; the minority has the right to remain in society and try to become the majority; but in the Church it is so easy to resort to church discipline and cast out those whom the majority does not like.

But, as said, it is dangerous for many reasons to establish bare majority rule in the Church. Perhaps the United Church is even now already beginning to taste the bitterness of the dangerous methods that were brought into use in order to displace Augsburgers from society and congregation.

For first, the Church is dependent upon God’s Word and its authority and truth before there is any talk of majority and minority. This limitation of majority rule is of such surpassing significance that where it is not heeded and respected, there the Church soon ceases to be Church. In other words, it is a fundamental matter in the Church that the Word be allowed to retain its proper authority.

Further, there is this circumstance to be taken into account, that there is no guarantee that the majority alone has good forces and sound views. At times there may be serviceable men and wise counsel in the minority. Therefore to exclude the minority from all influence and cast it out of the body, in order that it may never become the majority, is almost madness for those who would gather God’s Church and not scatter it.

And, not to go through everything, there is finally this peculiar feature of the Church’s apostolic and biblical order, that the congregation is free, self-standing, and independent in its very essence, so that it stands under no human or ecclesiastical authority. Congregations may cooperate with other Congregations in all freedom and brotherliness; but one congregation cannot rule over another and command what it shall do. There is not a hint of anything of that kind in Scripture. And therefore to arrange matters in such a way that one abolishes the congregation’s independence and responsibility is not scriptural.

For my part, it is in any case clear that a dominion of an ecclesiastical-political majority over the Congregations dissolves the fellowship of the Congregations instead of gathering it. Compulsion and tyranny will necessarily abolish the Church’s unity, however much they themselves lay claim to promoting it.

Free cooperation without compulsion is the proper ecclesial and Christian condition. Doubtless it is also that which leads to the happiest results. The so-called societies that have practiced forced cooperation have hitherto at any rate attained only divisions, because they have wished to carry on by force and majority what was meant to be free.

If therefore among our freedom-loving people, and among believing men who have responsibility before God, a real ecclesial life and brotherly cooperation are to come into being, then we must get out of the worldly and political tracks into which the societies have led us.

It will take a long time to get out of them, I know that well. And it is something to be thankful for, that one says this to us, that it goes no otherwise in the Free Church than elsewhere. Very well, what may one fairly expect? We must try to grow in age and wisdom.

A free cooperation is brought about only by the approval of all who cooperate. Every Congregation and every individual in the Congregation is asked whether they will be with us. It costs a great deal of labor, which forced cooperation seems to avoid. Yet it is done, because there is interest and spirit behind it. Therefore it also costs less money.

The short time in which we have tried free and brotherly cooperation has at any rate shown that, with all the difficulties which the coercive societies have diligently enumerated to us, it has nevertheless gone forward. And our conviction is this: if only we practice freedom and love in full earnest and with full integrity, then our cooperation will not be poor because it is free.

But if by freedom we mean only to be free from labor, and by love, that we will share no pain or toil with anyone, but only enjoy the pleasant company of others and other good things, then it is self-evident that we are good for nothing. Let us practice freedom and love, and voluntary cooperation shall show results that answer to the divine promises resting upon his free servants and self-sacrificing and self-denying children.

If all does not go on in sheer glory and joy, then let us not lose heart at once. Fair-weather Christians, who can come forward only when everything is compliant, never become truly free, because they have altogether too much to accommodate themselves to and to take account of. Let us go forward courageously, and let each bear his own burden, because it is right, and bear one another’s burdens, because thus we fulfil the Law of Christ.

If one will then call this fanaticism, it is nothing other than what has always been faith’s lot and portion. And yet faith is an equally sure and good way, by whatever name it may be called.

\subsection{Every free Congregation has, like each individual believer, through the prompting of God’s Spirit and by the right of love, the right to do good and to labor for the salvation of souls and the awakening of spiritual life as far as its ability and strength reach; in such free, spiritual activity it is restricted neither by parish ties nor by the boundaries of societies.}

“Christ’s love compels us” is the principle that lies behind and beneath all truly Christian action.

God loved; therefore he sent the Son into the world, that he should be the Savior of the world.

God’s Son loved; therefore he came into the world to seek and to save that which was lost. Out of free and living love he became the servant of the many and gave his life as a ransom-payment for them.

And God’s Spirit is he who pours out love into believers’ hearts, so that they are driven by the same mind as Jesus Christ himself. They love, because they are born of God and have his life in their hearts.

He who has the true faith and the living love cannot be still where there is a world round about that is led astray by falsehood, perishes through unbelief, and dies the eternal death in its sins. He must seek to “do likewise” as the merciful Samaritan, as Jesus Christ. The salvation of souls lies upon the heart of the saved, and they must try to get more to go with them on the way to eternal life. “Come and go with us” is the invitation of the saved to all who need salvation.

But what applies to the individual Christian must also apply to the Congregation. Or how should the Body of Christ have another will, prompting, and desire than Christ himself?

For the Congregation exists in order to manifest Christ’s life and be his instrument and servant. What else can it do, then, than do the works of him who has bought it with his blood and bestowed his Spirit upon it?

Christ’s love reaches to all; he will draw all to himself, he will gather all to his great supper. Human beings unto the uttermost ends of the earth are the object of his love.

The Congregation therefore cannot be without mission. It must farther forward and farther out. Its whole being is ordered toward the spread of God’s Kingdom; it has come forth by mission, and it is appointed to carry on mission. The one follows from the other.

And here mission is taken in its most general sense, embracing those who are near and those who are far away, both what is called Inner Mission and Jewish Mission and Heathen Mission. The meaning is to seek after that which is lost, wherever it is found.

Love from God is unrestricted by everything except God and God’s Word. When therefore men often employ human arrangements in order to shut God’s Word and Spirit out from themselves, love tries to leap over walls and break through barriers. And whether the barriers are called Chinese walls or Catholic boundary-bolts or state-church parish ties makes no difference to love.

Or how was it with the merciful Samaritan? Was that his “district,” where the poor traveler was assaulted between Jericho and Jerusalem? Was there first a question about the boundaries? Or was not mercy free according to its own eternal law and nature, so that it has mercy on whom it has mercy, without asking after anything else than the need and misery?

When therefore one has spread abroad among worldly church-people the delusion that one need only enter into a certain society, and then one is sheltered from lay preachers, awakening, and other pestilence, then it is the Congregation’s calling to overcome such false notions and break down such fences. Nowhere have souls the freedom and the right to sleep in peace in spiritual matters, so surely as the Congregation can do anything to awaken them and pull them out of a burning house. Even though the spiritually sleeping have hired pastor and sexton ever so much to protect their sleep and make it heavier and more secure, the Congregation is nevertheless equally bound to labor for the salvation of the sleeping. Jesus visited every synagogue into which he could enter, and the Temple itself became witness to the manifestation of his mighty power despite the displeasure of his adversaries. And Paul preached wherever he came, despite synagogue rulers. And the Congregation in Jerusalem wrote to all Congregations and admonished them.

There is only one consideration which the Congregation takes; it is love’s own consideration, namely this question: Can we accomplish something good? Where no benefit is done, love commands that one refrain from doing harm.

If the Congregation’s freedom and right were only ever so little acknowledged among us, then party spirit and its many thorn-fences between the hearts of the people would soon have to yield. The brotherhood, which rests upon fellowship with Christ in his Congregation, would then be able to thrive; and the brotherhood which exists only within the party would disappear before the true and rightful brotherhood.

But there is still a long way to the goal that must be reached if we are to be pleasing to God. From Congregation to Congregation the holy bonds of brotherhood must be knit without the societies’ shutting partition-walls.

Will it ever be reached? At all events, it is a worthy goal for which to labor and suffer.

Many are they who have suffered and suffer exceedingly much because of our people’s ecclesiastical sundering. There is only one remedy, and that is the working of God’s Spirit unto the liberation of hearts and Congregations. If the Congregation becomes what it ought to be, then party spirit receives its death-blow.

It is therefore extraordinarily characteristic to see the zeal with which the ecclesiastical-political organs among us work against the congregational idea. They seem to see only misfortune and ruin in it.

But that Congregation should have care for Congregation, that Congregation should have trust and confidence toward Congregation, that a believing person should be well received from Congregation to Congregation, that there should be an invisible chain linking those who serve the same Lord and contend for the same cause—how can that be any misfortune? Does this not accord with God’s Word and God’s Spirit, with the innermost essence of Christianity and with apostolic practice?

The matter is this, that men believe in the society which human beings build by constitutions; they esteem so lightly the fellowship which God builds by the Spirit.

Therefore we will conclude with the same with which we began: the task is to bring forth the Congregation from the rubble and to rebuild the ruins of the Lord’s temple among men. And if we do not reach far, which is indeed to be expected under our own great and dreadful sloth and weakness and the opponents’ enormous preponderance, humanly speaking, then well and good: we count it a high thing to labor for the right goal, and we are altogether certain of this, that the day shall come when all truth-loving souls shall see it clearly in the Spirit’s light, that no ecclesial labor was more urgently required in the present time than labor for the Congregation.

It is the only labor that guards against the danger of Hierarchy on the one side and the danger of Individualism on the other. Let it go forward with power.

